\documentclass[11pt,a4paper]{article}
\usepackage[T1]{fontenc}
\usepackage[utf8]{inputenc}
\usepackage[margin=23mm]{geometry}
\usepackage{amsmath,amssymb,mathtools,booktabs,longtable,array}
\usepackage{xurl}
\usepackage[colorlinks=true,linkcolor=blue,urlcolor=blue,citecolor=blue]{hyperref}
\hypersetup{pdftitle={PV and Battery Trading: Current Mathematical Model},pdfauthor={bess pv trading repository}}
\usepackage{enumitem}
\setlist{itemsep=3pt,topsep=5pt}
\setlength{\parindent}{0pt}
\setlength{\parskip}{6pt}
\setlength{\emergencystretch}{3em}
\newcommand{\E}{\mathbb{E}}
\newcommand{\R}{\mathbb{R}}
\newcommand{\pos}[1]{\left[#1\right]_+}
\DeclareMathOperator{\clip}{clip}
\title{PV and Battery Trading in MGP and MI Auctions\\Current Mathematical Model and Execution Workflow}
\author{\texttt{bess\_pv\_trading} repository}
\date{Implementation reviewed on 20 October 2026}
\begin{document}
\maketitle
\begin{abstract}
This document describes the implemented optimization and execution system for a
single photovoltaic farm and an AC-coupled battery sharing one point of
interconnection. It distinguishes auction plans, immutable market commitments,
physical feedback execution, and ex-post settlement. The shared Gurobi formulation
supports a risk-neutral scenario tree. The current historical January experiment
instead uses one deterministic rolling weather forecast and perfect auction-price
information, followed by execution against measured AC production. Separate
full-information PV+BESS and PV-only optimizations quantify an ex-post operating
value of storage. Equations below follow the current source code; they are not a
proposal or a claim of a fully causal historical trading strategy.
\end{abstract}
\tableofcontents
\newpage

\section{Scope and implementation map}

One farm has PV generation, an AC-coupled BESS, and a common AC grid connection.
There is no local electricity demand. PV can be exported, stored, or curtailed;
the battery can charge from PV or grid imports and discharge toward the grid.
The market sequence is MGP (called DA in code), MI-A1, MI-A2, and MI-A3 (IDA1,
IDA2, IDA3). Every submitted signed quantity is assumed to execute in full at the
auction clearing price. Bid prices, acceptance uncertainty, market impact, and
continuous intraday trading are outside this formulation.

The implementation has three distinct objects:
\begin{enumerate}
\item An optimization plan, which includes the current actionable auction order,
future auction recourse, and a battery/PV schedule over the remaining horizon.
\item An immutable commitment ledger, updated at submission and publication.
\item A physical trajectory produced by a feedback/protection policy using
realized interval-average PV. Its energy, wear, and deviations are settled
separately from the optimization objective.
\end{enumerate}

\begin{longtable}{>{\raggedright\arraybackslash}p{0.40\textwidth}>{\raggedright\arraybackslash}p{0.53\textwidth}}
\toprule Source & Responsibility \\
\midrule\endhead
\path{src/bess_pv_trading/domain.py} & Validated grids, assets, scenarios, information nodes, states and solver contracts.\\
\path{src/bess_pv_trading/optimization/model.py} & Physical, trading, product-resolution, optional imbalance constraints and objective.\\
\path{src/bess_pv_trading/optimization/degradation.py} & Throughput coefficient and segmented cycle-depth model.\\
\path{src/bess_pv_trading/optimization/validation.py} & Numerical validation independent of the solver constraint objects.\\
\path{src/bess_pv_trading/simulation/synthetic.py} & Shared event replay and auction submission/publication accounting.\\
\path{src/bess_pv_trading/simulation/replay.py} & Residual scenario conditioning and bounded physical execution.\\
\path{src/bess_pv_trading/simulation/gme_oracle.py} & Manual-data historical price oracle, rolling PV and full-information asset comparisons.\\
\path{src/bess_pv_trading/simulation/settlement.py} & Exact-grid actual settlement and expected/realized KPI reconciliation.\\
\path{src/bess_pv_trading/forecasting/pv.py} & Weather-to-AC PV conversion.\\
\path{src/bess_pv_trading/markets.py} & Configurable local auction gates, publication events and DST-aware delivery grids.\\
\bottomrule
\end{longtable}

\section{Time, units, signs and information}

\subsection{Delivery intervals and auction products}
A local delivery day is partitioned into physical intervals
$\mathcal T=\{0,\ldots,T-1\}$, with starts $\tau_t$, ends $\tau_{t+1}$ and durations
\begin{equation}
 \Delta_t=(\tau_{t+1}-\tau_t)/\text{one hour}.
\end{equation}
The historical study uses quarter-hours, so $\Delta_t=0.25$. A local day is
constructed between successive midnights in \texttt{Europe/Rome}, then stored in
UTC: $T$ is 92, 96, or 100 for spring-DST, ordinary, or autumn-DST days. Repeated
local hours are distinguished by their actual UTC offset/fold. No equation
assumes that every day has exactly 24 elapsed hours.

A physical interval is not necessarily an independent market product. Before
1 October 2025, DA products are hourly; the MI auctions in the current 2025
history are already quarter-hourly. Physical battery controls are quarter-hourly
at all four optimization stages. For the four equal-duration quarters of one
native DA hour $h$, the implementation enforces
\begin{equation}\label{eq:hourly}
 q_{\mathrm{DA},n,t}=q_{\mathrm{DA},n,t'}\quad(t,t'\in h).
\end{equation}
These are equal MWh shares of one hourly order, whose total is
$\sum_{t\in h}q_{\mathrm{DA},n,t}$. The hourly clearing price is repeated as an
objective coefficient on those shares. The quarters do not become separate DA
trades or separate observed prices. There is no corresponding date-based
hourly equality on battery charging or discharging. Native product resolution
and optional physical-control grouping are separate configuration fields.

The current study's configured local calendar is:
\begin{center}
\begin{tabular}{llll}
\toprule Auction & Gate & Results time used & Eligible delivery \\
\midrule
DA & D$-1$ 12:00 & D$-1$ 13:30 & Whole day D\\
IDA1 & D$-1$ 15:00 & D$-1$ 15:30 & Whole day D\\
IDA2 & D$-1$ 22:00 & D$-1$ 22:30 & Whole day D\\
IDA3 & D 10:00 & D 10:30 & D 12:00 onward\\
\bottomrule
\end{tabular}
\end{center}
Submissions are planned five minutes before the gate in the supplied study.
These times, date validity, result-time provenance and cancellations are explicit
YAML settings. In particular, the historical DA publication time is a study
assumption. Gate time and the start of IDA3 delivery are different events.
Recorded publication timestamps can override configured times.

An auction $a$ has an eligibility set $\mathcal T_a$. A documented/configured
cancellation removes that auction from the active set. Missing data alone does
not establish a cancellation: unexplained missing or incomplete curves fail
validation. Absent auctions create no trade variables and no submission solve.

\subsection{Notation and electrical basis}
\begin{longtable}{p{0.20\textwidth}p{0.17\textwidth}p{0.54\textwidth}}
\toprule Symbol & Unit & Meaning \\
\midrule\endhead
$E$ & MWh & Nominal internal stored-energy capacity.\\
$E^{\min},E^{\max}$ & MWh & Operational energy limits derived from SOC percentages.\\
$\overline P$ & MW AC & One symmetric battery charging/discharging rating.\\
$P^{\rm imp},P^{\rm exp}$ & MW AC & Directional shared POI ratings.\\
$\eta_c,\eta_d$ & fraction & Complete AC-bus-to-storage and storage-to-AC efficiencies.\\
$\sigma$ & fraction/hour & Configured hourly self-discharge fraction.\\
$V_{s,t}$ & MW AC & Forecast available PV in scenario $s$.\\
$\widetilde V_t$ & MW AC & Effective realized PV used for execution.\\
$c_{n,t},d_{n,t}$ & MW AC & Battery charge/discharge commands at information node $n$.\\
$u_{n,t}$ & MW AC & PV curtailment.\\
$i_{s,t}$ & MW AC & Uncurtailed PV injection into the common AC bus.\\
$g_{s,t}$ & MW AC & Net grid exchange; positive is export, negative is import.\\
$e_{s,t}$ & MWh & Internal battery energy at an interval boundary.\\
$q_{a,n,t}$ & MWh & Signed auction adjustment; positive sale, negative purchase.\\
$Q_{s,t}$ & MWh & Final cumulative contracted delivery position.\\
$\pi_{a,s,t}$ & EUR/MWh & Auction-price coefficient for interval $t$.\\
$w_s$ & fraction & Scenario probability, summing to one.\\
\bottomrule
\end{longtable}

Power variables are AC at the common bus/POI. Battery energy is internal stored
energy. Nameplate PV capacity is DC; measured historical \texttt{power} is already
AC kW, converted once to MW. A separate PV-output/transformer conversion is used
only when predicting AC PV from modeled DC. BESS efficiency is not applied to
measured PV or direct PV exports.

\subsection{Forecast vintages and observability}
A weather forecast cycle $v$ has source time $t_v$ and availability time $A_v$.
Without explicit source publication times, the configured convention is
$A_v=t_v+4$ hours. At decision cutoff $k$, the rolling policy chooses the latest
eligible cycle with $A_v\le k$ and required remaining-horizon coverage. Weather
cycles about six hours apart are successive forecast vintages, not ensemble
members. The current dataset supplies one deterministic weather member.

Measured interval-average PV becomes available after interval end plus configured
telemetry delay. A forecast for future intervals is not replaced by future
realized power in the ordinary auction replay. If a decision occurs inside an
unfinished quarter-hour, the next boundary state is projected through that
interval with forecast PV and the same execution policy. Its final realized
average cannot yet be observed. The hindsight counterfactuals deliberately have
future PV and price information and are labeled as such.

Policy-treated observations carry evidence times. Plateau evidence can require
later samples; two-sided interpolation requires its right endpoint and applicable
mask evidence. A finalized corrected time series therefore differs from the
as-of view available to a past decision. This remains relevant even when a study
uses actual-price foresight.

\section{PV power from weather forecasts}

\subsection{Irradiance and orientation}
The input mappings declare forecast shortwave radiation as GHI and direct
radiation as horizontal beam irradiance (BHI), with the latter an explicit
working assumption. Hourly radiation means are end-labelled and redistributed
uniformly over contained quarters. Other fields follow their declared
instantaneous/accumulation semantics rather than the same blanket interpolation.
Solar geometry is evaluated at each quarter-hour midpoint using configured
coordinates and farm elevation.

For eligible irradiance data,
\begin{equation}
 G_{\rm DHI}=G_{\rm GHI}-G_{\rm BHI},\qquad
 G_{\rm DNI}=\frac{G_{\rm BHI}}{\cos\theta_z}
 \quad\text{away from the low-sun cutoff}.
\end{equation}
The code uses pvlib's DNI helper with configured low-sun handling and a finite
DNI cap or invalidation rule. BHI exceeding GHI is either invalidated or explicitly
clipped and flagged. Optional Erbs decomposition can replace missing/unusable
direct input under its declared conditions; it is not a general repair of
contradictory source values.

Fixed arrays use metadata tilt and a declared source-to-pvlib azimuth conversion.
Tracking IDs 1 and 3 use \texttt{pvlib.tracking.singleaxis} with configurable axis
tilt, axis azimuth, rotation limit, backtracking and ground-coverage ratio. Type
3 currently requires the explicit front-side-only approximation; rear-side
bifacial gain is omitted. The current geometry uses an acknowledged Milan
location proxy rather than verified farm coordinates.

With panel tilt $\beta$, incidence angle $\alpha$ and ground albedo $a$, the
isotropic transposition can be written
\begin{equation}
 G_{\rm POA}=G_{\rm DNI}\pos{\cos\alpha}
   +G_{\rm DHI}\frac{1+\cos\beta}{2}
   +aG_{\rm GHI}\frac{1-\cos\beta}{2}.
\end{equation}

\subsection{Temperature, DC output and AC conversion}
Wind speed is the magnitude of the configured horizontal wind components,
$v=\sqrt{v_u^2+v_v^2}$. The Faiman temperature model is
\begin{equation}
 T_{\rm mod}=T_{\rm air}+\frac{G_{\rm POA}}{U_0+U_1 v}.
\end{equation}
The supplied wind height is explicitly acknowledged as a module-wind proxy.
The resulting temperature is passed to pvlib's PVWatts DC calculation:
\begin{equation}
 P_{\rm DC}=\max\left\{0,\ P_{\rm DC,STC}
 \frac{G_{\rm POA}}{1000}
 \left[1+\gamma(T_{\rm mod}-T_{\rm ref})\right]\right\}.
\end{equation}
DC power is scaled to AC by the configured explicit PV-output efficiency
$\eta_{\rm PV}$, or by an explicitly fitted DC-to-AC coefficient $\kappa$:
\begin{equation}
 \widehat P_{\rm AC}=
 \begin{cases}
 \eta_{\rm tr}\eta_{\rm PV}P_{\rm DC},&\text{explicit-efficiency route},\\
 \eta_{\rm tr}\kappa P_{\rm DC},&\text{fitted-coefficient route}.
 \end{cases}
\end{equation}
A configured AC cap clips this output if present; the night mask and validity
flags are then applied. The present study uses an unfitted explicit output factor
and a configurable separate transformer factor. This is not a full calibrated
module-and-inverter electrical model. Temperature and wind affect production
through temperature; irradiance and orientation affect POA. Cloud/dew-point
fields retained in the store are not additional direct PV predictors here.

Invalid required inputs remain invalid unless an explicit policy applies.
Historical GME studies may use their configured night-zero forecast policy and
a bounded fallback for missing realized-power intervals. Both the number and
identity of substituted intervals are recorded. All counterfactuals use the same
\emph{effective} realized PV series, including these substitutions.

\section{Optimization at an auction decision}

Each residual solve uses only the remaining suffix of the delivery grid. In the
equations below that suffix is reindexed from zero, with $T$ denoting its final
boundary; its original UTC product identifiers are retained for commitments.

\subsection{Scenario tree and nonanticipativity}
At decision $k$, let $\mathcal S_k$ be represented leaf scenarios and let
$\mathcal A_k$ be the open active auctions. An information tree assigns an auction
bid owner $n_a(s)$ to each auction/leaf and a control owner $n_t(s)$ to each
interval/leaf. Curtailment may have a separate owner $m_t(s)$ under the selected
execution-information convention.

Decisions are node-owned variables, not independent leaf decisions followed by a
penalty for disagreement. For example,
\begin{equation}
 c_{s,t}=c_{n_t(s),t},\quad d_{s,t}=d_{n_t(s),t},\quad
 q_{a,s,t}=q_{a,n_a(s),t},\quad u_{s,t}=u_{m_t(s),t}.
\end{equation}
Thus two scenarios with the same observable history share actionable decisions.
A latent PV-member label is not available information by itself. Constraints
apply to every represented leaf, including zero-probability leaves if supplied;
probabilities only weight the objective. In the current one-scenario historical
price oracle, the tree degenerates to deterministic planning with known future
prices and a forecast PV trajectory.

\subsection{Physical variables and balances}
For every scenario and remaining interval:
\begin{align}
 i_{s,t}+u_{s,t}&=V_{s,t}, & i_{s,t}&\ge0, & 0\le u_{s,t}&\le V_{s,t},\label{eq:pvbalance}\\
 g_{s,t}&=i_{s,t}+d_{s,t}-c_{s,t},\label{eq:bus}\\
 -P^{\rm imp}\le g_{s,t}&\le P^{\rm exp}.\label{eq:poi}
\end{align}
Curtailment appears only in the PV balance. It cannot erase bought grid energy
or battery charging energy. A negative grid exchange must be absorbed by net
battery charging since there is no local demand. The model does not explicitly
assign charging power to its PV/grid source; that partition is an accounting
convention used in post-processing.

With $z_{n,t}\in\{0,1\}$ denoting charging mode,
\begin{equation}\label{eq:power}
 0\le c_{n,t}\le\overline P z_{n,t},\qquad
 0\le d_{n,t}\le\overline P(1-z_{n,t}).
\end{equation}
This excludes simultaneous charge and discharge, uses one symmetric AC rating,
and introduces no separate battery-inverter power limit.

SOC percentage parameters $s_{\min},s_{\max}$ define
\begin{equation}
 E^{\min}=E\frac{s_{\min}}{100},\qquad
 E^{\max}=E\frac{s_{\max}}{100},\qquad
 E^{\min}\le e_{s,t}\le E^{\max}.
\end{equation}
For $\rho_t=(1-\sigma)^{\Delta_t}$, the transition is
\begin{equation}\label{eq:soc}
 e_{s,t+1}=\rho_t e_{s,t}
    +\eta_c c_{s,t}\Delta_t-\frac{d_{s,t}\Delta_t}{\eta_d}.
\end{equation}
One AC MWh charged stores $\eta_c$ MWh. Delivering one AC MWh consumes
$1/\eta_d$ internal MWh. In combined mode the two configured efficiencies each
cover their entire directional path. In converter-and-cell mode,
$\eta_c=\eta_{c,\rm conv}\eta_{c,\rm cell}$ and
$\eta_d=\eta_{d,\rm conv}\eta_{d,\rm cell}$. The alternatives cannot be mixed.
The round-trip factor is $\eta_c\eta_d$. Self-discharge acts on stored energy,
not on a power variable, and is separate from cycle degradation cost.

The initial boundary is fixed to measured/carried/projected state $e^k$:
\begin{equation}
 e_{s,0}=e^k.
\end{equation}
If the terminal target is configured, also impose
\begin{equation}\label{eq:terminal}
 e_{s,T}=e^{\rm target}.
\end{equation}
In the GME adapter, initial/terminal percentages are converted using nominal
capacity. A \texttt{null} terminal SOC omits~\eqref{eq:terminal}. Prepared and
synthetic boundary-state contracts use explicitly named MWh values. Neither
version adds a terminal energy salvage value to the objective.

\subsection{Auction quantities and physical delivery}
Let $Q^{\rm fix}_t$ contain all previously submitted commitments, priced or still
pending. The new final position on leaf $s$ is
\begin{equation}\label{eq:position}
 Q_{s,t}=Q^{\rm fix}_t+
       \sum_{a\in\mathcal A_k:\,t\in\mathcal T_a}q_{a,n_a(s),t}.
\end{equation}
An ineligible auction/interval has no trade variable. The model allows signed
sales/purchases; they are adjustments to prior positions, not gross generation.
For every eligible trade,
\begin{equation}\label{eq:tradepoi}
 -P^{\rm imp}\Delta_t\le q_{a,n,t}\le P^{\rm exp}\Delta_t.
\end{equation}
For every open auction $a$ in chronological order, define the cumulative planned
position after that auction as
\begin{equation}
 C^a_{s,t}=Q^{\rm fix}_t+
 \sum_{b\in\mathcal A_k:\,b\le a,\,t\in\mathcal T_b}q_{b,n_b(s),t}.
\end{equation}
It satisfies
\begin{equation}\label{eq:cumulative}
 -P^{\rm imp}\Delta_t\le C^a_{s,t}\le P^{\rm exp}\Delta_t,
 \qquad -P^{\rm imp}\Delta_t\le Q_{s,t}\le P^{\rm exp}\Delta_t.
\end{equation}
There is no independent commercial volume cap. The directional POI ratings apply
to individual auction adjustments, cumulative positions after each auction, and
physical delivery. An offsetting later trade is allowed, subject to its own
POI-derived limit and the resulting position limit. Thus the model permits
bounded cross-auction spread trading even without a battery; its revenues can be
present in both asset counterfactuals.

With \texttt{trading.optimize\_imbalances: false}, which is the current farm-910
policy, planned delivery is exact in every interval:
\begin{equation}\label{eq:exact}
 g_{s,t}\Delta_t=Q_{s,t}.
\end{equation}
No positive/negative imbalance variables or imbalance direction binary are
created. Forecast deviations are not an additional dispatch resource.
An optional configured product-group delivery equality can additionally impose
$\sum_{t\in h}(g_{s,t}\Delta_t-Q_{s,t})=0$. It is redundant with~\eqref{eq:exact}
and does not replace it in the disabled-imbalance mode.

\subsection{Current objective and sunk cash flows}
For throughput degradation with coefficient $c_{\rm thr}$ and an optional import
fee $f$, the current residual objective is
\begin{equation}\label{eq:objective}
 \max J_k=\sum_{s\in\mathcal S_k}w_s\left[
 \sum_{a\in\mathcal A_k}\sum_{t\in\mathcal T_a}\pi_{a,s,t}q_{a,n_a(s),t}
 -c_{\rm thr}\sum_t(c_{s,t}+d_{s,t})\Delta_t
 -f\sum_t\pos{-g_{s,t}}\Delta_t\right].
\end{equation}
The implementation uses an exact import-power linearization with a direction
binary only when $f$ is nonzero. The farm-910 study sets $f=0$.

Past auction cash flow and already incurred degradation are sunk at residual
reoptimization. Priced commitments remain in the external ledger, and pending
unpriced submitted quantities enter $Q^{\rm fix}$ without adding an optimizable
cash-flow term. Their future clearing value is constant with respect to all
remaining decisions and is excluded from $J_k$. This avoids double counting.
The initial pre-DA solve has no prior commitments; its objective spans the
modeled whole-day opportunity set. Residual objective values therefore should
not be summed across auction solves to compute daily profit.

There is no explicit curtailment penalty, risk measure, capital cost, fixed O\&M,
tax, brokerage fee, ramp constraint, minimum dwell time, or calendar-aging charge
in~\eqref{eq:objective}. Curtailment can affect revenue through delivery choices,
but is not independently priced. The horizon can exploit inventory when its
terminal target is free; this matters when comparing asset value.

\subsection{What is free at each of the four decisions}
\begin{center}\small
\begin{tabular}{p{0.13\textwidth}p{0.37\textwidth}p{0.39\textwidth}}
\toprule Decision & New free auction quantities & Previously submitted fixed quantities\\
\midrule
DA & DA, IDA1, IDA2, IDA3 plans & None\\
IDA1 & IDA1, IDA2, IDA3 plans & DA\\
IDA2 & IDA2, IDA3 plans & DA and IDA1\\
IDA3 & IDA3 eligible products & DA, IDA1 and IDA2\\
\bottomrule
\end{tabular}
\end{center}
At each decision, remaining-horizon charge/discharge modes, curtailment, grid
powers and energy states are free subject to their balances, boundary state,
nonanticipativity, product rules and terminal condition. Executed past intervals
are absent from the residual grid. Future-auction quantities are plans until
their own submission. Only the current auction order is committed from that
solve. Before noon, IDA3 cannot change orders for the morning; physical control
replanning of remaining earlier intervals must respect their already fixed
market positions.

\section{Battery degradation models}

\subsection{Current AC throughput approximation}
For one trajectory,
\begin{equation}
 H=\sum_t(c_t+d_t)\Delta_t,\qquad D_{\rm thr}=c_{\rm thr}H.
\end{equation}
Both directional AC throughputs are charged once. The YAML can directly supply
$c_{\rm thr}$ or derive it from replacement economics and reference cycle life.
Let $p_{\rm USD}$ be pack USD/kWh, $x$ be USD per EUR, $p_{\rm extra}$ additional
EUR/kWh replacement cost, $s_R$ salvage fraction, $N$ reference cycle life and
$d_{\rm ref}$ reference internal depth of discharge. Then
\begin{align}
 R&=1000E\left(\frac{p_{\rm USD}}{x}+p_{\rm extra}\right)(1-s_R),\label{eq:replacement}\\
 H_{\rm life}&=N d_{\rm ref}E\left(\frac{1}{\eta_c}+\eta_d\right),\\
 c_{\rm thr}&=\frac{R}{H_{\rm life}}.\label{eq:cthr}
\end{align}
One reference internal cycle charges $d_{\rm ref}E/\eta_c$ AC MWh and discharges
$d_{\rm ref}E\eta_d$ AC MWh, explaining the normalization. There is no additional
factor of two or half applied to~\eqref{eq:cthr}. The configured end-of-life
capacity fraction is descriptive metadata, not an extra multiplier or a dynamic
capacity update.

The current inherited farm-910 YAML uses a historical generic pack-price proxy,
a dated FX conversion, and an illustrative reference cycle life. Those are study
inputs with source metadata, not an estimate of a specific installed BESS cost.
The resolved $R$ and $c_{\rm thr}$ are saved. The cost is financial wear; it does
not reduce $E$, $\overline P$, or efficiencies as the simulation proceeds.

\subsection{Optional segmented depth-dependent aging}
The optional \texttt{xu\_segmented} mode follows the configured segmented
cycle-depth construction described in the repository's Xu notes and the author
paper~\cite{xu}. It is distinct from the current throughput study.
For a configurable convex stress curve $\Phi(d)=\alpha d^\beta$, $\beta\ge1$,
partition nominal internal capacity into $J$ equal segments of size $E/J$.
The marginal AC discharge cost of segment $j=1,\ldots,J$ is
\begin{equation}
 m_j=\frac{RJ}{\eta_d E}
       \left[\Phi(j/J)-\Phi((j-1)/J)\right].
\end{equation}
It is nondecreasing with depth. The degradation charge is
\begin{equation}
 D_{\rm Xu}=\sum_t\sum_{j=1}^J m_j d_{j,t}\Delta_t.
\end{equation}
Only discharging is charged under this convention. A discharging half-cycle
incurs the convention's depth cost, while a charging half-cycle does not; do not
add an extra factor of one half. Nonzero throughput cost is rejected in this mode.

Segment flows and energies obey
\begin{align}
 \sum_j c_{j,t}&=c_t,& \sum_j d_{j,t}&=d_t,&\sum_j h_{j,t}&=e_t,\\
 h_{j,t+1}&=h_{j,t}+\Delta_t(\eta_c c_{j,t}-d_{j,t}/\eta_d),&
 0\le h_{j,t}&\le E/J.
\end{align}
Segment energies represent cycle-depth memory, not fixed absolute SOC buckets.
Charging and discharging both use shallow available segments first. For each
$j<J$, the implementation adds priorities $v^c_{j,t},v^d_{j,t}\in\{0,1\}$:
\begin{align}
 \sum_{\ell>j}c_{\ell,t}&\le\overline P v^c_{j,t},&
 h_{j,t+1}&\ge(E/J)v^c_{j,t},\\
 \sum_{\ell>j}d_{\ell,t}&\le\overline P v^d_{j,t},&
 h_{j,t+1}&\le(E/J)(1-v^d_{j,t}).
\end{align}
Using a deeper segment therefore requires the shallower segment to end full
when charging or empty when discharging. This chooses a deterministic shallow-
first representative, including otherwise indifferent charging allocations.
Segment actions have the same information ownership as aggregate battery power.

At a declared fresh study start, initialize
\begin{equation}
 h_{j,0}=\clip\big(e_0-(j-1)E/J,\ 0,\ E/J\big).
\end{equation}
Afterwards the segment vector must be carried with actual execution; an arbitrary
residual solve cannot reconstruct it from aggregate SOC. Previous throughput or
aging requires explicit memory. Xu mode currently rejects nonzero self-discharge.
It remains a finite-segment approximation and omits calendar, temperature and
current-rate aging, capacity-fade feedback and unknown pre-study cycle history.

\section{Submission, publication and reoptimization workflow}

Before each active auction submission, the GME replay refreshes the latest
available PV forecast over unexecuted intervals, projects a boundary if necessary,
and solves the residual planning problem. Previously submitted quantities,
including still-unpriced ones, are fixed. At submission, only the current
auction's node-shared action is extracted. Later planned auction quantities are
not booked and can change.

At publication, known prices are updated, and actual contractual revenue
$\pi^{\rm real}_{a,t}q^{\rm sub}_{a,t}$ is entered exactly once in the ledger.
Publication does not reopen a market. The current GME experiment forces
\texttt{reoptimize\_physical\_controls=false}; it uses four auction solves for
four active auctions, not an extra set of publication solves or 96 physical
optimizations. Physical feedback is algebraic. Synthetic replay can separately
enable interval-control reoptimization, so its solve count need not be four.

At the next auction, closed commitments are carried, remaining open opportunities
are reconsidered, and past executed state/aging remain fixed. In the actual-price
oracle all future auction prices are already supplied as a benchmark, so price
publication does not remove price uncertainty there. In a forecast-based route,
known results condition subsequent price distributions without changing the
commitments that produced them.

\section{Physical feedback with realized PV}

\subsection{Feasible interval power range}
The current policy is \texttt{track\_committed\_grid}. For interval $t$ its target is
\begin{equation}
 g_t^\star=Q_t^{\rm committed}/\Delta_t.
\end{equation}
The controller uses effective realized interval-average PV $\widetilde V_t$;
this is an idealized interval simulator, not a subsecond telemetry/control model.
It does not disclose this final average to the earlier auction optimizer.
Given actual initial energy $e_t$ and retained energy $\bar e_t=\rho_t e_t$,
the feasible charging/discharging maxima are
\begin{align}
 C_t&=\min\left\{\overline P,
           \frac{\pos{E^{\max}-\bar e_t}}{\eta_c\Delta_t}\right\},\\
 D_t&=\min\left\{\overline P,
           \frac{\eta_d\pos{\bar e_t-E^{\min}}}{\Delta_t}\right\}.
\end{align}
First clip the planned charge/discharge commands to these maxima, obtaining
$c_t^0,d_t^0$. Simultaneous positive requested charging and discharging after
protection is rejected. Initially curtail only available PV:
$u_t^0=\min\{\widetilde V_t,\pos{u_t^{\rm plan}}\}$.

\subsection{Grid-tracking policy}
Release scheduled curtailment if it can help meet a deficit:
\begin{align}
 g_t^0&=\widetilde V_t-u_t^0+d_t^0-c_t^0,\\
 \bar u_t&=u_t^0-\min\{u_t^0,\pos{g_t^\star-g_t^0}\}.
\end{align}
The desired net battery dispatch and its feasible projection are
\begin{align}
 b_t^{\rm req}&=g_t^\star-\widetilde V_t+\bar u_t,\\
 b_t&=\clip(b_t^{\rm req},-C_t,D_t),\\
 c_t&=\pos{-b_t},\qquad d_t=\pos{b_t}.
\end{align}
Finally curtail any remaining surplus that battery adjustment cannot absorb:
\begin{equation}
 u_t=\bar u_t+
 \min\{\widetilde V_t-\bar u_t,
             \pos{\widetilde V_t-\bar u_t+b_t-g_t^\star}\}.
\end{equation}
Thus excess PV first decreases discharge, then increases charging as needed and
feasible, then is curtailed. A PV shortfall first releases planned curtailment,
then reduces charging or increases discharge within the actual SOC box. The
remaining tracking error is preserved if battery power/energy cannot supply it.

POI protection follows. If net import exceeds its rating, charging is reduced.
If export exceeds its rating, discharge is reduced first and remaining excess
PV is curtailed. Grid power is recomputed consistently with the adjusted actions.
Actual energy then advances using~\eqref{eq:soc} with these executed powers, and
bounds are checked before accepting the new state. Both throughput/Xu memory and
cumulative financial wear are advanced from executed actions, not from the old
auction plan. Protection is mandatory under either execution policy; the
alternative \texttt{follow\_battery\_plan} skips grid-tracking adjustment.

There is no controller constraint forcing the actual terminal energy to the
auction plan's target. Forecast errors and grid tracking can cause an ending-SOC
difference even if every optimization imposed~\eqref{eq:terminal}. This is visible
in the current terminal-50\% study and is not equivalent to guaranteeing 50\%
realized SOC. Any hard actual-terminal policy would be an additional controller
design and could create different imbalances/cash flows.

\subsection{PV source allocation for reporting}
For the executed trajectory, use the explicit PV-first accounting convention
\begin{align}
 P^{\rm PV\to B}_t&=\min\{\widetilde V_t-u_t,c_t\},\\
 P^{\rm PV\to G}_t&=\widetilde V_t-u_t-P^{\rm PV\to B}_t,\\
 P^{\rm PV\to curt}_t&=u_t.
\end{align}
They sum to $\widetilde V_t$ in every interval. Grid-supplied charging and battery
discharge are not counted as PV production. These are post-processing flows;
there are no separate source-attribution decision variables in the MILP.

\section{Settlement and expected-versus-realized accounting}

Actual net deviation is retained at the native quarter-hour:
\begin{equation}
 \delta_t^{\rm real}=g_t^{\rm real}\Delta_t-Q_t^{\rm committed}.
\end{equation}
Positive means more net export than promised. With the selected single Terna
NORD settlement price $\lambda_t$, the modeled cash flow is
\begin{equation}
 S_t=\lambda_t\delta_t^{\rm real}.
\end{equation}
A positive price therefore does not mean every imbalance is penalized: positive
excess export is credited and a negative deviation is debited. Negative prices
are allowed and reverse these cash-flow signs. The general settlement helper
also accepts two directional curves:
\begin{equation}
 S_t=\lambda_t^+\pos{\delta_t^{\rm real}}
       -\lambda_t^-\pos{-\delta_t^{\rm real}}.
\end{equation}
This documents the implemented study accounting, not every possible Italian
participant-specific settlement item or regulatory regime. The workbook adapter
selects the declared NORD macrozone and the imbalance-price column, retaining
other imported fields as provenance. Product starts and ends must match the
execution grid exactly. Prices are neither hourly-averaged nor filled across
gaps; missing enabled settlement prices are an error.

Actual daily auction cash flow and settled operating profit are
\begin{align}
 A^{\rm real}&=\sum_{a,t}\pi^{\rm real}_{a,t}q^{\rm sub}_{a,t},\\
 J^{\rm real}&=A^{\rm real}+\sum_t S_t
       -D^{\rm executed}-f\sum_t\pos{-g_t^{\rm real}}\Delta_t.
\end{align}
A model deterrence penalty is not an actual settlement charge and is not
subtracted a second time by actual post-processing. If settlement is unavailable,
profit before settlement is reported and fully settled profit remains unknown.

Let $J^{\rm init}$ and $A^{\rm init}$ be the initial auction solve's expected
profit and auction revenue. Saved mismatches are
\begin{equation}
 \Delta J=J^{\rm real}-J^{\rm init},\qquad
 \Delta A=A^{\rm real}-A^{\rm init},\qquad
 r_J=100\frac{\Delta J}{|J^{\rm init}|}.
\end{equation}
The percentage is undefined/\texttt{null} when its expectation is zero; unsettled
profit comparisons remain unknown. Differences include later replanning and
physical execution effects, not only price-forecast error.

The replay also evaluates each stored auction snapshot as an open-loop diagnostic
against realized PV and the same settlement prices. Its planned battery commands
are protected physically. Deviations are measured against that snapshot's planned
grid schedule (which can include assumed future trades) or the final contract,
as labeled in the recorded comparisons. It is not another auction
optimization. Because IDA3 covers a shorter residual horizon, a common-horizon
comparison is also saved; unaligned whole-horizon totals should not be compared
as if their delivery intervals were identical.

\section{Perfect-information schedules and the PV-only baseline}

\subsection{PV+BESS versus the feedback controller}
After replay, one additional full-day problem uses $V_t=\widetilde V_t$ and all
actual auction-price curves, with a single scenario and all auction orders
jointly free subject to their native products, eligibility, POI and physical
constraints. It retains the BESS ratings, efficiencies, SOC box, degradation,
initial energy and configured terminal target. Call its achieved solver
objective $J_{\rm PV+BESS}^{\rm PI}$.

This is deliberately a hindsight plan rather than a causally executable policy.
It is compared with the feedback controller's actual dispatch, curtailment and
SOC, and its objective is saved as the solver's \texttt{ObjVal}. Four missing
measurement intervals, if permitted by the selected run, are handled through
that run's explicit fallback; diagnostics record their count and identities.
The schedule figure shows both trajectories and their battery-power difference
through the final day boundary. It does not assert that feedback should reproduce
the hindsight trajectory or terminal energy.

\subsection{PV-only asset counterfactual}
When \texttt{experiment\_outputs.pv\_only\_baseline: true}, a second independent
solve uses the same optimizer builder with
\begin{equation}
 E=0,\quad\overline P=0,\quad E^{\min}=E^{\max}=0,\quad e_0=e_T=0,
 \quad D=0.
\end{equation}
Battery powers and energy variables are consequently fixed to zero; the code
reuses one physical/trading formulation instead of implementing a second,
potentially inconsistent PV-only model. Unused variables can be removed by
solver presolve. PV injection/curtailment and auction trades remain free subject
to the same effective realized PV, native products, active auctions and POI.
Call its achieved objective $J_{\rm PV}^{\rm PI}$.

The incremental operating-value KPI is
\begin{equation}\label{eq:value}
 \Delta J_{\rm BESS}^{\rm PI}
       =J_{\rm PV+BESS}^{\rm PI}-J_{\rm PV}^{\rm PI}.
\end{equation}
The BESS side includes its own modeled degradation; the PV-only side has no
battery wear. Both sides retain permitted cross-auction spread trading, so
market-trading revenue common to both is not credited to battery ownership.
The KPI does not subtract realized feedback wear from a different ideal schedule.
Each counterfactual reports its own objective components. Directly subtracting
a PV-only hindsight objective from the realized PV+BESS feedback profit would
mix information and execution policies and is not the asset-value KPI.

This comparison is one-day operating value before capital expenditure, fixed
operating costs, financing and taxes. A free terminal target can also use initial
stored energy or leave different inventory; the selected 50\% initial/terminal
hindsight comparison avoids that particular inventory difference for the BESS
model. Realized feedback may still finish at a different SOC. No residual-energy
valuation, annualization, payback, NPV or investment optimization is implemented.
Those need a representative horizon and additional cost/cash-flow assumptions.

Setting the switch to \texttt{false} skips the PV-only solve and its model, JSON,
CSV and figure. Incremental-value metrics become \texttt{null} and an explicit
baseline-enabled KPI is false. The perfect-information PV+BESS-versus-feedback
comparison still runs. Four active auctions therefore produce four replay
solves plus one PV+BESS hindsight solve and, when enabled, one PV-only solve.
Counterfactual solve diagnostics are separate from replay solve counts.

\section{Forecast prices and the general stochastic route}

The historical GME study described above has no price noise and uses actual
prices in every auction decision. The repository separately implements
chronological point-price studies and an artifact-to-optimization route for
uncertain prices and probabilistic PV members.

A generic selected linear model is $\widehat\pi=x^\top\widehat\beta$, possibly
modeling a change from yesterday's price rather than the absolute level.
Calendar, lagged prices, daily profiles and configured regional load/renewable
forecasts form features. OLS, ridge, weighted ridge, lasso, elastic net and Huber
estimators share training-only preprocessing and portable fitted artifacts.
Validation selects specifications before later calibration and evaluation.

For an auction's native delivery vector, Gaussian calibration uses out-of-sample
errors to estimate a bias $b_a$ and temporal covariance $\Sigma_a$. Paths are
\begin{equation}
 \pi_a=\mu_a(\mathcal H_a)+b_a+\varepsilon_a,
 \qquad\varepsilon_a\sim\mathcal N(0,\Sigma_a).
\end{equation}
An example covariance policy is
\begin{equation}
 \Sigma_a=(1-\gamma)S_a+\gamma\operatorname{diag}(S_a)+\epsilon I,
\end{equation}
with saved shrinkage and variance-floor settings. Bias and zero-mean innovations
are distinct; finite Monte Carlo samples need not have exactly zero sample mean.
IDA conditional means may be reconstructed from a predecessor-price curve plus
a learned spread, using only ancestor sampled/known clearing curves at their
information nodes. Innovations are temporally correlated within an auction
equation and independent across auction equations; their conditional means
propagate inter-auction dependence. Negative simulated prices are not clipped.

For $B$ branches at each of $K$ unknown auction results there are up to $B^K$
price paths per supplied PV member, before history merging and configured limits.
Two branches at four unknown auctions give up to 16 price paths per member.
Known results are fixed, not sampled. The prepared route forms joint scenario
mass from PV-member and price-path masses; price draws shared across members
represent the declared independent product-measure assumption rather than
estimated PV/price dependence. Member identity remains latent until observations
distinguish histories. The actual-price oracle repeats a common price curve
across members and adds no artificial price noise.

Hourly and quarter-hour price regimes and DST templates remain distinct.
Historical retrieved revisions lacking publication evidence are marked
retrospective. A model fit cutoff, static feature availability, forecast cycle,
and source timestamp must satisfy the selected decision cutoff policy.
Explicit retrospective opt-in does not establish an operationally causal
backtest. The prepared examples need their complete supplied tables/artifacts;
they are not a fully populated alternative historical strategy by default.

\section{Optional optimized-imbalance formulation}

This section describes supported code that is disabled in the current farm-910
study. When \texttt{optimize\_imbalances=true}, replace~\eqref{eq:exact} by
\begin{align}
 \delta^+_{s,t}-\delta^-_{s,t}&=g_{s,t}\Delta_t-Q_{s,t},\\
 0\le\delta^+_{s,t}&\le M_t y_{s,t},\\
 0\le\delta^-_{s,t}&\le M_t(1-y_{s,t}),\qquad y_{s,t}\in\{0,1\}.
\end{align}
Here $M_t=(P^{\rm imp}+P^{\rm exp})\Delta_t$, optionally reduced by a configured
imbalance-energy cap. The common $M_t$ is valid because physical and contractual
net energy use the same directional POI range. The sign binary prevents
simultaneous positive/negative deviations, including with negative prices.
Add to each scenario's objective
\begin{equation}
 \sum_t\left[\lambda^+_{s,t}\delta^+_{s,t}
      -\lambda^-_{s,t}\delta^-_{s,t}
      -\kappa(\delta^+_{s,t}+\delta^-_{s,t})\right].
\end{equation}
If this mode has zero settlement price, zero penalty and no useful delivery
restriction, an optimized contract need not be backed by delivered energy.
That is why the present experiment disables optimized deviations and computes
physical mismatches only afterwards. The optional formulation is a mathematical
study mode, not an assertion that deliberate imbalance strategies are permitted.

\section{Data provenance, outputs and reproducibility}

Farm inputs are local read-only Zarr/Parquet stores. API and workbook records use
\path{data/api/market.duckdb}. Raw/normalized snapshots, source identities,
checksums, revision and retrieval metadata support repeated offline studies.
Historical publication remains a separate field/convention. Provider acquisition
is an explicit user-run command; neither the MILP builder nor the historical
experiment downloads missing inputs. DuckDB writers and analysis processes must
be coordinated; complete cached days can be studied before an entire acquisition
window is complete.

YAML controls study date, calendar validity, source mappings, PV and data policies,
BESS and POI characteristics, state, degradation, solver and outputs. One include
is resolved relative to the containing YAML, with recursive mapping overrides.
The current 50\% study is configured by
\path{configs/full_experiments/gme_oracle_2025-01-01_farm910_perfect_pv_comparison_terminal_soc_50.yaml}.
Its inherited assumptions include a 5 MWh battery, 2.5 MW symmetric power,
10--90\% SOC box, 0.97 combined efficiencies and a 6.638 MW directional POI.
These are reproducible study settings, not confirmed anonymous-farm hardware.

Each historical run is grouped under
\path{outputs/full_experiments/gme_oracle/<run>/farm_<id>/<day>/} with
\texttt{plots/}, \texttt{csv/}, \texttt{json/}, and optional \texttt{models/}.
The effective configuration, input provenance and artifact manifest accompany
auction plans, immutable ledger, actual dispatch, used PV vintages, settlement
and objective comparisons. The main KPI distinctions are:
\begin{itemize}
\item \texttt{total\_gain\_eur}: actual feedback profit after available settlement;
\item \texttt{perfect\_information\_objective\_eur}: achieved full-information PV+BESS objective;
\item \texttt{pv\_only\_perfect\_information\_objective\_eur}: achieved optional PV-only objective;
\item \texttt{battery\_incremental\_operating\_value\_eur}: the difference in~\eqref{eq:value};
\item initial expectation, realized-minus-expected revenue/profit, and separate executed versus counterfactual degradation components.
\end{itemize}
Figures preserve interval widths and both available boundary points, including
next-day local midnight. Held hourly/quarter-hour prices and powers do not create
new samples; SOC uses the actual boundary states. Snapshot settlement diagnostics
include their different delivery windows and common-horizon comparisons.

Solver settings use relative MIP gap $10^{-4}$ (0.01\%) by default, with no absolute
MIP gap configured. Objective comparisons use achieved \texttt{ObjVal}. Status,
relative gap, model sizes and timing remain recorded. Before a schedule is used,
independent validation recomputes physical balances, limits, native product
consistency, information sharing, degradation and objective components. Execution
separately checks actual SOC and power/POI protection. Software consistency does
not resolve uncertain farm metadata, forecast bias or historical publication
vintages.

\section{Current limitations and extension boundaries}

The current empirical workflow is a single-day perfect-price benchmark with one
weather member, approximate geometry and an idealized interval-average controller.
Measured AC is effective observed production, not an independently estimated
unconstrained irradiance-to-generation potential; hidden plant curtailment/faults
in source production are not identified automatically. Configured cleaning and
fallback choices remain part of the experiment.

The model currently omits continuous trading, order acceptance, market impact,
multiple-farm network coupling, ancillary services, intrainterval dynamics,
ramping/startup constraints, dynamic capacity loss, calendar aging, thermal
battery dynamics, fixed operating costs and investment cash-flow optimization.
The generic \texttt{backtest} CLI remains synthetic-only; the separate historical
\texttt{experiment-gme-oracle} workflow is the implemented manual-data replay.
Future uncertain-price evaluations need their own causal inputs and comparison
protocol rather than relabeling the current oracle results.

The original implementation-status report is retired. This source and its PDF
are the maintained mathematical description; run manifests and saved solver
results are the numerical evidence for particular configurations. Practical
commands and configuration examples are in the root \path{README.md}.

\begin{thebibliography}{9}
\bibitem{xu} B. Xu, J. Zhao, T. Zheng, E. Litvinov and D. S. Kirschen,
\emph{Factoring the Cycle Aging Cost of Batteries Participating in Electricity
Markets}. Author preprint and published paper links retained by the repository:
\url{https://arxiv.org/abs/1707.04567} and
\url{https://doi.org/10.1109/TPWRS.2017.2733339}.
\bibitem{pvlib} pvlib Python documentation. The installed library implements the
solar-position, single-axis tracking, isotropic transposition, Faiman and PVWatts
functions called by \path{forecasting/pv.py}.
\url{https://pvlib-python.readthedocs.io/}.
\bibitem{gme} GME market timing references and period-specific study provenance
are recorded in \path{markets.py} and the selected study's YAML. The January
historical configuration explicitly labels assumed result-publication clocks.
\bibitem{code} Repository source files listed in Section 1 and the effective
configuration saved with each run are authoritative for implemented behavior.
Conceptual references do not replace inspection of those inputs and policies.
\end{thebibliography}
\end{document}
